\chapter{Finite Fields and Arithmetic Directions}
\label{ch:finite-fields}

The theory does not stop over \(\C\): it also moves to finite fields and number
fields.  This is possible because the chord--tangent formulas are rational in
the coordinates.  Once the cubic equation and the addition formulas are written
with coefficients in a field \(k\), the same formulas define a group \(E(k)\)
whenever the cubic is smooth.

\begin{theorem}[Hasse bound]
\label{thm:hasse}
For an elliptic curve \(E/\F_p\),
\[
  \#E(\F_p)=p+1-a_p,\qquad |a_p|\le 2\sqrt p.
\]
\end{theorem}

\begin{proof}[Status]
This is a deep theorem, not proved in this textbook.  It is included to mark the
transition from complex geometry to arithmetic geometry.
\end{proof}

Cryptography uses the finite abelian group \(E(\F_p)\).  Adding points is fast
because it is given by explicit rational formulas.  Reversing repeated addition,
recovering \(n\) from \(P\) and \(nP\), is the elliptic-curve discrete logarithm
problem.

The Tate module is another arithmetic object.  For a prime \(\ell\ne p\), the
\(\ell^n\)-torsion groups fit into an inverse system, and
\[
  T_\ell E=\varprojlim E[\ell^n]\cong \Z_\ell^2.
\]
If \(E\) is defined over \(\Q\), the absolute Galois group acts on these torsion
points:
\[
  \rho_\ell:\Gal(\overline{\Q}/\Q)\to \GL_2(\Z_\ell).
\]

\section*{Exercises}
\addcontentsline{toc}{section}{Exercises}

\exerciseintent{Count finite-field points by a concrete residue table.}
\begin{exercise}
Count the points on \(y^2=x^3+2x+1\) over \(\F_5\).  Make a table for
\(x=0,1,2,3,4\), compute \(x^3+2x+1\pmod 5\), compare with the square set
\(\{0,1,4\}\), and then add the point at infinity.
\end{exercise}

\exerciseintent{Identify the field operations used by the rational addition formulas.}
\begin{exercise}
In the chord--tangent formulas, list the field operations that appear
\(+, -, \times,\) and division by a nonzero denominator.  From that list, state
the calculation rule that lets the same formulas define addition over any field
where the cubic is smooth, and name the exceptional denominator cases that must
be handled separately.
\end{exercise}

\exerciseintent{Place a finite-field point count inside Hasse's interval.}
\begin{exercise}
Write the Hasse bound as the explicit interval
\[
  p+1-2\sqrt p\le \#E(\F_p)\le p+1+2\sqrt p.
\]
Then interpret the interval geometrically.
\end{exercise}

\exerciseintent{Separate algebraic ingredients valid over general fields from complex-analytic ones.}
\begin{exercise}
Make a two-column table.  In the first column, list algebraic ingredients that
still make sense over a general field: the cubic equation, \(O\), divisors, and
rational chord--tangent formulas.  In the second column, list analytic
ingredients that depend on \(\C\): lattices, \(\C/\Lambda\), \(\wp\), periods,
and complex Lie groups.
\end{exercise}
